\section{Introduction}
A purchaser of an operating commitment buys a response to incentives, not an externally fixed stream of service. A transfer can induce participation without determining how long an operator remains active, whether the operator surrenders, or which of several privately optimal policies is selected. The difficulty is amplified when continuation values are supplied by an approximate dynamic solver. A small value error need not identify the derivative of value, and an accurately evaluated selected policy need not identify the purchaser's worst response. This paper develops Neural Bellman Operators as an interface between dynamic approximation, response information, and verifiable contract choice.

The analysis begins with a simple observation: a fixed policy's payoff is affine in a change in its reward, even when the optimized value is nonsmooth and the policy responds discontinuously. Finite value messages therefore restrict response moments jointly. We characterize the sharp set relative to these messages and distinguish it from the response set of a fully known dynamic model. In that model, occupancy flow equations and one global private-value budget characterize all near-optimal responses. Lagrange duality then expresses support of the response set through directed reward queries. A neural operator can propose those queries or policies, but an independent value enclosure, rather than the training loss or a chosen tie-breaking rule, determines the economic statement.

This distinction gives a practical information-design result. If the purchaser values duration and fee receipts together, a joint reward perturbation can be substantially more informative than separate coordinate secants. If transition structure is available, a further whole-interval response graph discards actions that cannot be optimal at any common fee in a cell. The graph does not silently replace a common contract by state-specific fees: it is an explicit outer relaxation whose remaining excess is charged in the upper bound. Both constructions admit independently checked certificates and record the cost of the information used.

We apply the resulting bounds to participation, enforcement capacity, and a continuous surrender fee. An executable contract must provide sufficient compensation and adequate service under every admitted response. A competing contract is allowed its most favorable response in the upper comparison. Covering every fee interval then gives a global purchaser-regret bound without treating a dense menu as a continuum. An explicit behavioral convention separates exact offers with attained responses from positive-tolerance offers when an initial financial mandate is open. It prevents a closure optimizer from being offered as an unavailable action.

The economic contribution concerns the substitution between continuation opportunities and enforcement. We characterize the least initial-state fee that removes a valuable surrender option by a payoff-to-surrender ratio. We also characterize the uniform enforcement requirement by the obstacle gap above the no-surrender value. Enlarging adjustment opportunities raises the latter value state by state and weakly lowers the uniform requirement. This result allows endogenous consumption, financial positions, repeated adjustment, and stopping; it is not confined to an irreversible-upgrade model. The initial-state comparison requires an additional calculation because changing reachability can change which states matter for exact implementation.

In the recorded stochastic settlement economy, allowing a continuous fee removes the compulsory-term difference found on the original coarse menu: both regimes admit a certified term-one offer. That result is retained. The new enforcement comparison explains what survives the menu refinement. Adjustment raises the no-surrender private value by more than $0.03545$ and lowers the initial-state no-surrender threshold by more than $0.07209$. Structural response-graph certificates, rather than a rounded surrender probability, establish the zero-surrender side of this comparison. Continuous institutional counterfactuals change fee receipts, capacity curvature, and efficiency. A separate adverse parameter map preserves cases in which adjustment reduces the financial-mandate premium; the continuation-value result is not used to erase that distinction.

The computational argument separates approximation, policy evaluation, improvement, and certification. We retain the unfavorable original neural experiment and add three-seed comparisons across four action--horizon families, with nearest-anchor, polynomial, evaluated policy-bank, and exact dynamic-program controls. These experiments do not establish neural state-dimension superiority. Their role is to measure the actual interface: which proposals yield useful value intervals, what they cost, and which claims remain valid when a proposal fails. For a continuous diffusion, a direct Bellman-barrier theorem supplies the corresponding derivative, action-supremum, boundary, and stopping requirements. The stored-array certificates and the numerical diffusion-to-array error are distinct objects; a small arithmetic residual does not identify the latter.

\paragraph{Related literature.}
The use of inequalities to delimit counterfactual behavior connects the paper to revealed preference and inverse optimization. \citet{afriat1967} reconstructs preferences from expenditure choices; \citet{ahuja2001} studies optimization primitives consistent with prescribed solutions. Our messages instead concern counterfactual dynamic values, and the unknown object of immediate interest is a response moment. \citet{magnac2002} studies identification in dynamic discrete decisions; that econometric problem uses observed choice information and restrictions on primitives, rather than purchaser-designed value queries. The distinction between a known transition model, a class of observationally admissible models, and an oracle-only model class is essential to interpreting sharpness here. Robust convex optimization, as in \citet{bental1998}, supplies a natural language for using an outer set without pretending that every point in it is dynamically attainable.

The paper also builds on verification in computational economics. \citet{santos2005} studies the accuracy of simulations in stochastic dynamic models, \citet{kubler2005} distinguishes approximate from exact dynamic equilibria, and \citet{kubler2011} develops sufficient verification conditions. \citet{judd2017} gives lower bounds on approximation error; a lower error diagnostic must not be confused with the upper error enclosure needed for a purchaser certificate. Our decision-directed budget records signed query-error exposure and complements, rather than replaces, model approximation analysis. The standard numerical and approximation background includes \citet{judd1998}, \citet{stokey1989}, and \citet{barles1991}.

Neural stochastic control and operator-learning approaches include \citet{han2016}, \citet{han2018}, \citet{lee2024}, \citet{kim2025}, \citet{kim2026}, and \citet{cai2025}. Derivative approximation and physics-informed training are discussed by \citet{hornik1990} and \citet{raissi2019}. The proposal/evaluation distinction is related to policy iteration, policy switching, and information relaxation; see \citet{jacka2015}, \citet{chang2021}, and \citet{brown2010}. Reward transfer through cached policies and successor features is not new; see \citet{barreto2017}, \citet{nemecek2021}, and \citet{alegre2022}. Parametric-model verification likewise has a developed interval-cover literature, including \citet{quatmann2016}, \citet{junges2019}, and \citet{heck2025}. The present results use these ideas to make endogenous response, open-mandate attainability, and continuous institutional choice part of one checked economic comparison.

The next section states the economic environment. We then characterize response information, derive procurement and continuous-fee certificates, and analyze enforcement substitution and financial exposure. The last two sections connect Neural Bellman Operators to the certificates and distinguish fixed-target verification from continuous-model verification. Proofs and complete checking conventions are in the supplement. The preserved compendium contains the earlier operator, recursive-utility, dynamic-game, and numerical developments without abridgment.
